\section{What the adaptive curve retains and loses}\label{sec:structure}
The separation in Proposition~\ref{prop:separation} raises quantitative
questions. How far can a curve-only relaxation underestimate menu regret?
When is it exact? Can additional globally optimal presets compensate for
its failure? We answer these questions using the slopes available on the
active spectral faces. The convex-geometric arguments apply more generally
to compact families of affine functions; the spectral model supplies an
explicit face test and finiteness of the exceptional prices.

\subsection{Three feasible families}
Identify an affine line with its coefficient pair and let
$\mathcal L=\{(\alpha(P),c(P)):P\in\Gr\}$, a compact subset of $\R^2$.
Let $\mathcal G\subseteq\Gr$ contain every preset that attains $F(p)$ at
some $p\in[0,B]$, including every choice at ties. This set is compact:
it is the projection of the closed optimality relation in
$\Gr\times[0,B]$. Write $E_K^{\mathcal G}$ for the menu optimum restricted
to $\mathcal G$, and $\overline E_K$ for the optimum when each preset line
can be any member of $\operatorname{conv}\mathcal L$. The latter is a
mathematical relaxation; a convex combination of effective lines need not
be realizable by a projector. The lower envelope remains $F$. Compactness
ensures attainment of all three optima, and
\begin{equation}\label{eq:threefamilies}
 \overline E_K\le E_K\le E_K^{\mathcal G}.
\end{equation}

The geometric statements below hold for any nonempty compact coefficient
set $\mathcal L\subset\R^2$, with presets interpreted as its elements and
$F(p)=\min_{(\alpha,c)\in\mathcal L}(\alpha+pc)$. Convexity of
$\mathcal L$ is not assumed. The spectral content is the realization of
the sharp examples, the explicit active-face formula, and finiteness of
the exceptional set. In particular, these results do not assert a new
general scalarization or convex-duality principle.

For $a<b$ define the curve-only endpoint value
\begin{equation}\label{eq:curveoracle}
 D(a,b)=\max_{p\in[a,b]}
 \{F(p)-F(a)-(p-a)s\},\qquad s=\frac{F(b)-F(a)}{b-a}.
\end{equation}
For an interior price define the compact set of available optimal slopes
\begin{equation}\label{eq:slopeset}
 S(p)=\{c(P):J(P,p)=F(p)\}.
\end{equation}
The right and left derivatives of the compact affine envelope are
$F'_+(p)=\min S(p)$ and $F'_-(p)=\max S(p)$, respectively. One way to
see the right-derivative identity is to compare $F(p+h)$ with each
minimizing line at $p$, then take a convergent subsequence of minimizing
lines at $p+h$ as $h\downarrow0$. The reverse inequality follows from
optimality at $p$. The left identity follows in the same way. Thus the
superdifferential is $[\min S(p),\max S(p)]$.

\begin{lemma}\label{lem:relaxation}
$D(a,b)$ is the exact endpoint minimax value for
$\operatorname{conv}\mathcal L$. An attaining relaxed line is either a
globally optimal preset line or a convex combination of two preset lines
that are globally optimal at the same interior price.
\end{lemma}
\begin{proof}
Every relaxed line lies above $F$. Its regret relative to the chord at
any price is a convex combination of its two endpoint regrets, giving
$D(a,b)$ as a lower bound. Maximize $F(p)-sp$ at $p_*$. If
$a<p_*<b$, then $s\in[\min S(p_*),\max S(p_*)]$. Choose two optimal
lines whose slopes bracket $s$ and mix them to have slope $s$. They agree
at $p_*$, so the mixed line has both endpoint regrets equal to $D(a,b)$.
If $p_*=a$, select a minimizing line at $a$ of smallest slope. The
right-derivative optimality condition makes its slope at most $s$, and
its regret at $b$ is at most its regret at $a$. If $p_*=b$, use the
largest minimizing slope and the left-derivative condition. The same
argument uses the appropriate one-sided derivatives at $0$ and $B$.
\end{proof}

The acceptable set of a relaxed line is again an interval. Consequently
Theorem~\ref{thm:greedy}, with $D$ in place of $e$, computes
$\overline E_K$ from $F$ alone in exact arithmetic. It computes a lower
bound on the realizable menu problem, with the following sharp comparison.
This is a statement about values and covering intervals: recovery of an
actual projector still requires the model's active-face information.

\begin{theorem}\label{thm:factor}
For every $K\ge1$,
\begin{equation}\label{eq:factor}
 \overline E_K\le E_K\le E_K^{\mathcal G}\le2\overline E_K.
\end{equation}
The constants in $E_K\le2\overline E_K$ and
$E_K^{\mathcal G}\le2E_K$ cannot be decreased uniformly over the spectral
models and budgets. Sharpness already occurs in the limit with $K=1$.
\end{theorem}
\begin{proof}
An optimal relaxed menu supplies a partition into at most $K$ segments,
each served by one relaxed line at tolerance $\overline E_K$. Such a
partition is obtained from its closed interval cover as in
Theorem~\ref{thm:greedy}. On a segment $[a,b]$ choose any globally
optimal preset at $q=(a+b)/2$. Its two nonnegative endpoint regrets sum to
\[
 2F(q)-F(a)-F(b)\le2\overline E_K,
\]
because the covering relaxed line lies above $F(q)$ and its endpoint
regrets are at most $\overline E_K$. Each endpoint regret is therefore
at most $2\overline E_K$, and convexity extends this bound to the segment.
One such preset per segment proves the last inequality. The other
inequalities are inclusions of feasible families. The next example proves
sharpness.
\end{proof}

\paragraph{A continuum of indistinguishable curves.}
Extend Model M by taking $C_1=2I$ and $C_2=\diag(\kappa,0)$ for
$\kappa>2$, retaining $A=\diag(1,0)$, $r=1$, $B=1$, and zero overheads.
For every $\kappa$,
\begin{equation}\label{eq:kappa}
 F(p)=\min(-1+2p,0),\qquad
 \overline E_1=\tfrac12,\quad
 E_1=1-\kappa^{-1},\quad E_1^{\mathcal G}=1.
\end{equation}
Indeed, with $t=u_1^2$ the endpoint regrets are $1-t$ and
$\min(2-t,(\kappa-1)t)$. For $t\le2/\kappa$ their maximum is minimized
at $t=1/\kappa$, with value $1-1/\kappa$. For $t\ge2/\kappa$ its
minimum is $1$. For $p>0$ the only globally optimal projectors are the
appropriate coordinate projectors, including both at $p=1/2$; at $p=0$
only the first coordinate is optimal. The relaxed value follows from
\eqref{eq:curveoracle}. Hence
\[
 \frac{E_1}{\overline E_1}=2-\frac2\kappa\longrightarrow2
 \quad(\kappa\to\infty),\qquad
 \frac{E_1^{\mathcal G}}{E_1}=\frac{\kappa}{\kappa-1}
 \longrightarrow2\quad(\kappa\downarrow2).
\]
Also $E_1^{\mathcal G}=2\overline E_1$ throughout the family. The same
curve realizes every value of $E_1$ strictly between $1/2$ and $1$.
Common positive definite shifts preserve the regrets. At $\kappa=2$
the optimal face at $p=1/2$ contains every projector, so the limiting
claim does not assert continuity of the restricted feasible family.

The factor-two argument is the elementary two-objective weighted-sum
bound applied to endpoint regrets. Weighted scalarization and limitations
of supported solutions have a substantial approximation literature
\cite{bazgan2022,helfrich2024}; we do not claim that principle as new. Here its role is
to quantify the information loss in Proposition~\ref{prop:separation},
with sharp examples inside the specified spectral model.

\subsection{An exact criterion on active slopes}
Call $p\in(0,B)$ a \emph{slope-gap price} if
$S(p)\ne[\min S(p),\max S(p)]$. This definition concerns missing slopes,
not merely nondifferentiability: a corner can have every intermediate
slope available from an optimal projector.

\begin{theorem}\label{thm:exactness}
The identity $e(a,b)=D(a,b)$ holds for every $0\le a<b\le B$ if and
only if there are no slope-gap prices in $(0,B)$. In that case
\begin{equation}\label{eq:exactmenus}
 E_K=E_K^{\mathcal G}=\overline E_K\qquad(K\ge1).
\end{equation}
\end{theorem}
\begin{proof}
If every $S(p)$ fills its interval, the interior construction in
Lemma~\ref{lem:relaxation} uses one actual globally optimal preset of
slope $s$. Its boundary construction already uses one preset. Thus
$e(a,b)=D(a,b)$ with a globally optimal realizing preset on every
segment. Applying this to an optimal relaxed cover proves
\eqref{eq:exactmenus}.

Conversely let $s$ be a missing slope at $p_*\in(0,B)$. Compactness of
$S(p_*)$ implies
$l:=F'_+(p_*)<s<u:=F'_-(p_*)$. Fix
$R>(u-s)/(s-l)$. For small $h>0$ set $a=p_*-h$. The chord slope
from $a$ to $p_*$ tends to $u>s$, whereas the chord slope from $a$ to
$p_*+Rh$ tends to $(u+Rl)/(1+R)<s$. By continuity there is
$b\in(p_*,p_*+Rh)$ for which the chord slope is exactly $s$.
Concavity and $l<s<u$ make $p_*$ the unique maximizer of $F(p)-sp$.

Suppose an actual preset attained endpoint error $D(a,b)$. At $p_*$
the positive weighted average of its endpoint regrets equals
$J(P,p_*)-F(p_*)+D(a,b)\ge D(a,b)$. Since each endpoint regret is
at most $D(a,b)$, equality forces both regrets to equal $D(a,b)$
and $J(P,p_*)=F(p_*)$. Equal endpoint regrets force $c(P)=s$,
contradicting $s\notin S(p_*)$. Attainment on the compact projector
set therefore gives $e(a,b)>D(a,b)$.
\end{proof}

\paragraph{Computing the slope set from spectral faces.}
At an interior price let $\mathcal A(p)=\{L:F_L(p)=F(p)\}$.
For each $L\in\mathcal A(p)$ use the cutoff-face construction in
Theorem~\ref{thm:endpoint} to obtain its minimum and maximum costs
$c_{L,-}(p)$ and $c_{L,+}(p)$. Then
\begin{equation}\label{eq:unionfaces}
 S(p)=\bigcup_{L\in\mathcal A(p)}
 [c_{L,-}(p),c_{L,+}(p)].
\end{equation}
Every intermediate cost is realized by the explicit rotations in that
theorem. Moreover, a minimizing projector of an active site must have
that site among its cheapest sites: since $p>0$, a strictly cheaper site
would yield an objective below $F(p)$. Conversely, a globally optimal
projector minimizes at a cheapest active site. This proves
\eqref{eq:unionfaces}. The restriction $p>0$ is essential to this
argument; at zero price a tied site's raw cost need not be the
projector's effective cost.

In particular, a single site has no slope-gap prices, even at eigenvalue
collisions. Differentiability of $F$ throughout $(0,B)$ is another
sufficient condition, since then every $S(p)$ is a singleton. With
multiple sites, one tests whether the union in \eqref{eq:unionfaces}
fills its interval. A jump in an optimizing subspace alone does not
imply failure of the curve relaxation.

\begin{remark}
The equivalence concerns relaxation exactness on \emph{every subinterval}.
It is not a converse criterion for \eqref{eq:exactmenus} on a single
fixed domain, or for the optimality of supported presets alone. For
Model M on $[1/2-h,1/2+h]$, $0<h<1/4$, direct endpoint calculation
gives $e=e^{\mathcal G}=2h$ but $D=h$. Here $e^{\mathcal G}$ restricts
the endpoint problem to globally optimal presets. Thus that restriction
can be exact even when convexification strictly improves the objective.
\end{remark}

\subsection{A finite cost in additional presets}
\begin{theorem}\label{thm:repair}
If there are $q<\infty$ slope-gap prices in $(0,B)$, then
\begin{equation}\label{eq:repair}
 E_{K+\min(K,q)}^{\mathcal G}\le\overline E_K\le E_K
 \qquad(K\ge1).
\end{equation}
Without finiteness, $E_{2K}^{\mathcal G}\le\overline E_K$ still holds.
For the finite-dimensional spectral model, $q$ is always finite.
\end{theorem}
\begin{proof}
Partition an optimal relaxed cover into at most $K$ segments and replace
its line on each segment by an endpoint optimizer from
Lemma~\ref{lem:relaxation}. If a globally optimal preset realizes that
line, retain it. Otherwise its slope lies in a gap of $S(p_*)$ at an
interior maximizing price. Choose two globally optimal lines at $p_*$
whose slopes bracket the required slope. Their lower envelope is no
greater than their convex combination at every price. The two actual
presets therefore cover this segment at its relaxed tolerance.

Each replacement that needs two presets can be charged to a slope-gap
price strictly inside its segment. Distinct segment interiors are
disjoint, so each price is charged at most once. There are at most
$\min(K,q)$ extra presets. This is selection between two deterministic
presets after observing the price, not a randomized projector.

For finiteness, apply Rellich's analytic eigenvalue parameterization to
each affine symmetric pencil on an open interval containing $[0,B]$
\cite[Theorem~1.1]{barbarino2022}. Each $F_L$ is the minimum of the
finitely many analytic functions obtained by summing any $r$ of those
eigenvalue branches and adding $pd_L$. Consequently $F$ is the minimum
of a finite family of analytic functions. Discard identical functions.
Every pair then has finitely many zeros of its difference on the compact
interval $[0,B]$. Outside their finite union, the minimizing analytic
function is locally fixed, so $F$ is differentiable. A slope gap
requires nondifferentiability by \eqref{eq:slopeset}, proving finiteness.
No polynomial bound on $q$, or efficient exact enumeration of it, is
asserted.
\end{proof}

Let $N_{\rm rel}(\epsilon)$, $N(\epsilon)$, and
$N_{\mathcal G}(\epsilon)$ be the least menu sizes at tolerance
$\epsilon>0$ in the three families. They are finite, and
\begin{equation}\label{eq:cardinalityrepair}
 N_{\rm rel}(\epsilon)\le N(\epsilon)\le N_{\mathcal G}(\epsilon)
 \le N_{\rm rel}(\epsilon)+\min\{N_{\rm rel}(\epsilon),q\}.
\end{equation}
Thus the extra supported-menu budget is bounded independently of the
requested accuracy for each fixed spectral instance. The bound counts
exceptional prices, rather than the number of missing slope intervals
at one price. It can be conservative: different repairs may reuse a
preset. Some overhead is necessary: the family in \eqref{eq:kappa}
has $q=1$ and $E_1^{\mathcal G}=1>\overline E_1=1/2$, whereas its
two coordinate presets give $E_2^{\mathcal G}=0$. Thus the extra preset
cannot be removed uniformly even for $q=K=1$. We do not assert sharpness
of the additive $q$ term for every $q$ and $K$.

For $K>q$, \eqref{eq:repair} also gives
\begin{equation}\label{eq:asymptoticsandwich}
 \overline E_K\le E_K\le E_K^{\mathcal G}\le\overline E_{K-q}.
\end{equation}
If a positive asymptotic law $\overline E_K\sim cK^{-\alpha}$ has been
established for a particular curve, then both $E_K$ and
$E_K^{\mathcal G}$ have that same leading law. This transfers existing
scalar approximation results, such as suitable one-sided smooth-curve
laws~\cite{gaubert2011}; it does not establish a new exponent or assert
those regularity assumptions for every spectral instance. Piecewise
affine curves can instead admit a finite zero-regret menu.

\paragraph{Relation to representative-set approximation.}
Minimum-size convex Pareto approximation already distinguishes feasible
representatives from convex combinations~\cite{diakonikolas2008}; the
detailed treatment in \cite[Sections~5.2--5.3]{diakonikolas2011} also
exhibits loss from restricting representatives to extreme supported points.
That phenomenon is not new here. Nor is replacement by supporting pieces:
\cite[Lemma~2]{magnanti2012} gives a bounded-cardinality tangent replacement
for multiplicative scalar approximation. Our two-line replacement is the
corresponding elementary exposed-face construction for additive regret.
The refinement is to charge a replacement only at an interior price with
an unavailable active slope; disjoint covering intervals charge each such
price at most once. Combining this count with \eqref{eq:unionfaces} and
spectral finiteness gives the instance-wise additive overhead. It does not
improve the computational-complexity guarantees of those predecessors.
Appendix~\ref{app:comparison} makes the approximation criteria and the
relation to robust policy-count bounds explicit.
